%--------------------------------------------------------------------------------
%
%
%--------------------------------------------------------------------------------
\pagebreak
\section*{CBR - Compare and Branch}
\addcontentsline{toc}{section}{CBR - Compare and Branch}

%--------------------------------------------------------------------------------
\begin{iPageDescEntry}{Syntax}

	\texttt{CBR.<cond> RegR, RegB, target}
    
	where cond is EQ, LT, NE or LE.
\end{iPageDescEntry}

%--------------------------------------------------------------------------------
\begin{iPageDescEntry}{Format}

The CBR instruction uses the signed immediate S15 format:
	
\begin{flushleft}
\drawInstrFormat
  	{0,1,3,5,6.5,8.5,16}
  	{\OpCodeGrpBR,\OpCodeCBR,Reg R, cond, Reg B, ofs }
\end{flushleft}
\end{iPageDescEntry}

%--------------------------------------------------------------------------------
\begin{iPageDescEntry}{Description}

The \texttt{CBR} instruction compares the contents of \texttt{RegB} with \texttt{RegR}.  
If the condition specified in the \texttt{cond} field is satisfied, execution continues at the current instruction plus the offset. Otherwise, execution proceeds with the next instruction.  
The registers themselves are not modified.

\end{iPageDescEntry}

%--------------------------------------------------------------------------------
\begin{iPageDescEntry}{Conditions}

The \texttt{cond} field encodes the branch condition as follows:

\begin{center}
\begin{tabular}{|c|c|l|}
\hline
\textbf{Value} & \textbf{Mnemonic} & \textbf{Branch if} \\ \hline
0 & \texttt{EQ} & \texttt{RegR == RegB} \\ \hline
1 & \texttt{LT} & \texttt{RegR < RegB}  \\ \hline
2 & \texttt{NE} & \texttt{RegR != RegB} \\ \hline
3 & \texttt{LE} & \texttt{RegR <= RegB} \\ \hline
\end{tabular}
\end{center}

\end{iPageDescEntry}

%--------------------------------------------------------------------------------
\begin{iPageDescOpEntry}{Operation}
\begin{lstlisting}[style=iPageOpStyle]
if ( cond( RegR, RegB )) IA <- IA + ofs;
else                     IA <- IA + 4;
\end{lstlisting}
\end{iPageDescOpEntry}

%--------------------------------------------------------------------------------
\begin{iPageDescEntry}{Exceptions}
	None.
\end{iPageDescEntry}

%--------------------------------------------------------------------------------
\begin{iPageDescEntry}{Notes}
	None. 
\end{iPageDescEntry}
